\section{BCM systems, CM ground states, and exact neutrality}
\label{sec:bcm-neutrality}

This section incorporates Follow-up~1. HP denotes the supplied version
of Ha--Paugam, \emph{Bost--Connes--Marcolli systems for Shimura varieties I},
arXiv:math/0507101. CMR denotes Connes--Marcolli--Ramachandran,
\emph{KMS states and complex multiplication}, arXiv:math/0501424.
LLN denotes Laca--Larsen--Neshveyev, \emph{On Bost--Connes type systems
for number fields}, arXiv:0710.3452. Y denotes Yalkinoglu,
\emph{On arithmetic models and functoriality of Bost--Connes systems},
arXiv:1105.5022, including Neshveyev's appendix.
The section numbers for HP below are those of the supplied source:
its Dirichlet construction is \S8, not \S6.

\subsection{Positive partition functions and character insertions}

\paragraph{The source statements.}
HP \S7.3, first and second lemmas, gives the number-field Hamiltonian
$H e_{\mathfrak a}=\log N\mathfrak a\,e_{\mathfrak a}$ and partition
function $\zeta_F$. HP \S8.2 explicitly supposes that the character
is trivial at infinity. Its second and third definitions introduce
\emph{twisted} trace and partition function, and its lemma identifies
that twisted partition function with $L_F(s,\chi)$. Its final remark
leaves construction for nontrivial infinite components unresolved.
Thus the displayed construction covers even rational Dirichlet
characters, and finite-order characters over imaginary quadratic
fields; its own text does not supply the odd rational construction.
The zeros of odd Dirichlet functions are nevertheless covered by A.

\begin{proposition}[AC-B01: the exact analytic and thermodynamic consequences]
\label{AC-B01}
For a finite-order Hecke character $\chi$ with conductor $\mathfrak f$,
put $D_\chi e_{\mathfrak a}=\chi(\mathfrak a)e_{\mathfrak a}$, with
the value zero on ideals not prime to $\mathfrak f$. Then
\[
 \operatorname{Tr}(D_\chi e^{-sH})=L_F(s,\chi),\qquad \Re s>1.
\]
On the conductor-restricted Hilbert space used for the twisted trace,
the \emph{ordinary} positive heat trace is instead
\[
 Z_{F,\mathfrak f}(s)=
 \zeta_F(s)\prod_{\mathfrak p\mid\mathfrak f}
                      (1-N\mathfrak p^{-s}).                 \tag{B1.1}
\]
For $F=\mathbb Q$, A puts the nontrivial zeros of the primitive
twisted function on $\Re s=1/2$. For abelian $F/\mathbb Q$, A puts
the nontrivial zeros of $\zeta_F$ on that line. Finite Euler factors
deleted in (B1.1) have zeros on $\Re s=0$, and do not alter this
critical-strip conclusion. Continuation of these functions is a
classical input; the additional conclusion from A is their zero
location, not a new construction of the continuation.

The positive Gibbs state for the ordinary number-field system, when
$\beta>1$, is
\[
 \Phi_{\beta,w}(a)=\zeta_F(\beta)^{-1}
                 \operatorname{Tr}(\pi_w(a)e^{-\beta H}).
\]
Its established phase diagram is unchanged: there is one KMS state
for $0<\beta\le1$ and the extremal states for $\beta>1$, and at
zero temperature, are parametrized by
$\operatorname{Gal}(F^{\rm ab}/F)$. This is LLN Theorem~2.1
(\texttt{mainthm}), unconditional for every number field.
Meromorphic continuation does not make the heat operator trace
class for $0<\beta\le1$.
\end{proposition}
\begin{proof}
In the ideal basis, both operators in the first trace are diagonal.
The sum of absolute values is at most $\zeta_F(\Re s)$, which
converges for $\Re s>1$. Taking the trace therefore gives precisely
the ideal Dirichlet series with its primitive bad-prime convention.
Removing the zero entries is equivalent, for this trace, to
restricting to ideals prime to $\mathfrak f$. Without the insertion
the same basis gives the Euler product (B1.1).

For $F=\mathbb Q$ the first zero assertion is A. For abelian $F$,
Artin factorization expresses $\zeta_F$ as the product of the
primitive Dirichlet functions of the characters of
$\operatorname{Gal}(F/\mathbb Q)$. Each is holomorphic in the open
critical strip, so their zeros, counted with multiplicity, are the
union of those of the factors. A applies to all factors. A zero
of $1-N\mathfrak p^{-s}$ has $|N\mathfrak p^{-s}|=1$, hence
$\Re s=0$.

The state and its phase diagram are the cited LLN theorem; CMR
Theorem~5.1 (\texttt{KMSCM}) gives the imaginary-quadratic instance,
including its arithmetic action. These theorems are used, not
claimed as new consequences of A. Finally the positive trace
$\sum_{\mathfrak a}N\mathfrak a^{-\beta}$ diverges for
$0<\beta\le1$; the conductor-restricted series has the same
abscissa and divergence, by its nonzero finite Euler factor at one
and its positive coefficients. An analytic continuation changes
none of the positive eigenvalues $N\mathfrak a^{-\beta}$ or their
sum. It therefore cannot provide a normalized Gibbs density there.
\end{proof}

\begin{proposition}[AC-B02: a twisted trace is generally not a state]
\label{AC-B02}
For a nontrivial finite-order $\chi$, normalizing
$\operatorname{Tr}(D_\chi\,\cdot\,e^{-\beta H})$ by
$L_F(\beta,\chi)$ does not, in general, define a positive state,
even when that denominator is nonzero. In particular this failure
occurs for the even primitive character $\chi_5$ and for a cubic
character whose ideal values include $1,\omega,\omega^2$.
Inserting $D_\chi$ does not change the eigenvalues $\log N\mathfrak a$
of $H$. Hence the nontrivial zeros of the twisted partition
function are not thereby eigenvalues of this Hamiltonian.
\end{proposition}
\begin{proof}
On the positive rank-one projection onto $e_{\mathfrak a}$, the
normalized functional takes value
\[
 \frac{\chi(\mathfrak a)N\mathfrak a^{-\beta}}
      {L_F(\beta,\chi)}.                                    \tag{B2.1}
\]
For $F=\mathbb Q$, $\chi_5(2)=-1$, while $\chi_5(1)=1$ and
$L(\beta,\chi_5)>0$ for $\beta>1$: each real Euler factor is
positive. Formula (B2.1) is negative at $\mathfrak a=(2)$.
For a cubic character, no fixed nonzero complex scalar can send
all three numbers $1,\omega,\omega^2$ to nonnegative reals.
Thus positivity already fails on diagonal projections of
$B(\ell^2)$, the domain on which HP defines the twisted trace.
Equivalently, a positive normal functional on this algebra must
have a positive density operator, whereas the density displayed
here has these different phases. This proves the stated general
failure without asserting that every restriction to every smaller
subalgebra must fail. The diagonal formulas for $H$ and $D_\chi$
show the final assertion directly. A separate zero spectral
realization, such as AC-T04, uses different operators and a
different Hilbert space.
\end{proof}

For arbitrary Shimura BCM pairs, HP \S6.2,
Proposition \texttt{prop-construction-KMS}, constructs extreme
type-I Gibbs states in the convergence range. The following
Question \texttt{classif-etats-KMS-general} asks whether all
low-temperature extreme states arise from the Shimura variety.
The abelian number-field factorization above supplies no proof
of that general classification. The actual positive algebra,
dynamics, and measure classification remain necessary data.
There is also a harmless typographical omission of $\beta$ in
the first displayed cyclic vector in that proposition's proof:
its coefficient must be $\det(\phi(h))^{-\beta/2}$, as the
next computation and the Gibbs formula themselves require.

\subsection{Constructing the cubic ground-state observable}

\begin{proposition}[AC-B03: the radical is an actual arithmetic observable]
\label{AC-B03}
Let $K=\mathbb Q(\omega)$, $\alpha=\sqrt[3]p$, and let
$\mathfrak f$ be the conductor of $E_p/K$. In the arithmetic
endomotive of Y, one can construct a bounded finite-level
observable $a$ and a boundary idempotent $e$ such that
\[
                         a^3=p e.
\]
On a ground state belonging to the identity unit stratum,
$a$ has value $\alpha$; its three conjugate ground-state values
are $\alpha,\omega\alpha,\omega^2\alpha$.
The CMR CM arithmetic algebra likewise has ground-state values
generating $E_p$ and indeed all of $K^{\rm ab}$. This gives an
arithmetic realization of the radical, without an estimate of
the H04.4 ramp.
\end{proposition}
\begin{proof}
Y's finite-level algebra, equation \texttt{arithalg} in \S6,
is $E_{\mathfrak f}=\prod_{\mathfrak d\mid\mathfrak f}
K_{\mathfrak f/\mathfrak d}$. The quotient of finite
Deligne--Ribet sets constructed in AC-K03 induces an injection
\[
 B_p=E_p\times K^{d(\mathfrak f)-1}\lhook\joinrel\longrightarrow
 E_{\mathfrak f}.
\]
Take $a=(\alpha,0,\ldots,0)$ and $e=(1,0,\ldots,0)$ in $B_p$.
In the function description, $a(gH)=g\alpha$ on the unit
stratum and is zero on every other stratum. This is a
$G_K$-equivariant function, hence precisely an element of that
finite algebra. Pullback embeds it in Y's arithmetic subalgebra;
it is bounded because the finite-level function has finitely
many values. Its cubic identity is componentwise. Ground states
are evaluations on the invertible boundary, by LLN Theorem~2.1(iv)
and Y \S1.2, subsection \texttt{kmsinfty}. Evaluation at the
identity and Galois equivariance give the claimed values.

For CMR, Theorem~5.1 and the proof in \S5 under
``KMS states at zero temperature and Galois action'' give
$\varphi_{\infty,L}(f)=f(L,L)$, equation \texttt{Kground},
and the modular-value identity $F_\tau=K^{\rm ab}$,
equation \texttt{FtauKab}. In fact the values of functions
regular at $\tau$ form a field: if $f(\tau)\ne0$, then
$1/f$ is regular at $\tau$ as well. Thus the generated field
statement realizes $\alpha$ as the value of a rational
expression in finitely many such functions, regular at
$\tau$. This is an existence conclusion from the CM theorem,
not an explicit formula for such a modular function and not
an unproved identification of two arithmetic subalgebras.
None of these identities contains the logarithmic-derivative
coefficients or a bound on their sparse arithmetic windows.
\end{proof}

CMR's paragraph immediately after Theorem~5.1 supplies a further
necessary distinction. Inclusion of a subgroupoid in the
$\mathrm{GL}_2$ groupoid, even compatible with time evolution,
does not transport its KMS states to those of the ambient
system. Its explicit example is a finite full pair groupoid
inside an infinite one, with trivial dynamics: the first has
a normalized trace, the compact-operator algebra of the
second has none. Accordingly, CM ground states cannot be
transported to an ambient analytic zero theorem merely by
this inclusion.

\subsection{Neutrality at the units, at three, and at infinity}

\begin{proposition}[AC-B04: the exact finite and infinite alternatives]
\label{AC-B04}
For a unitary Hecke character of $K=\mathbb Q(\omega)$ with
infinite component
$\chi_\infty(z)=(z/|z|)^{-k}|z|^{2it}$, one has
\[
             \prod_{v<\infty}\chi_v(\epsilon)=\epsilon^k
                         \qquad(\epsilon\in\mu_6).          \tag{B4.1}
\]
A finite-order character necessarily has $k=t=0$.
Let $\mathfrak p=(\pi)\nmid3$ and let $r_\pi$ be the
nontrivial cubic residue character on
$(\mathcal O_K/\mathfrak p)^\times$, with values in $\mu_3$.
Then
\[
       r_\pi(\omega)=\omega^{(N\mathfrak p-1)/3}.            \tag{B4.2}
\]
If this value is $1$, $r_\pi$ descends to a finite-order ray
character of conductor $\mathfrak p$. Otherwise it can be
corrected by a uniquely specified cubic character modulo
$(3)=(1-\omega)^2$, yielding conductor exactly
$(3)\mathfrak p$. Correcting the same nontrivial unit colour
by an infinity type $k\equiv2$ or $4\pmod6$ instead yields
an infinite-order Hecke character.
\end{proposition}
\begin{proof}
A Hecke character is trivial on diagonal $K^\times$.
Applying this to $\epsilon\in\mu_6$, whose absolute value
is one, proves (B4.1). A continuous finite-image character
of the connected group $\mathbb C^\times$ is trivial.
The angular character and the radial character in the
displayed infinity type are both trivial only if $k=t=0$.
This also proves that a nonzero angular compensation cannot
remain finite order. For $\mathbb Q$, the same argument
at $-1$ gives $\chi_\infty(-1)=\chi_f(-1)$, since a sign
is its own inverse.

Formula (B4.2) is the definition of the cubic residue
character, applied to $\omega$. Also $r_\pi(-1)=1$:
its order divides both two and three. Since $K$ has class
number one, the ray class group for a finite modulus is
the group of residue units modulo the image of $\mu_6$.
Thus $r_\pi$ descends exactly when it is trivial on $\mu_6$,
equivalently when (B4.2) is one, or
$N\mathfrak p\equiv1\pmod9$.

The reduction map
\[
                 \mu_6\longrightarrow(\mathcal O_K/(3))^\times
\]
is an isomorphism. Its target has six elements. The images
of $-1$ and $\omega$ have orders two and three respectively:
$\omega-1$ is divisible by $1-\omega$ but not by its square.
They therefore generate six distinct units. Define $c$ on
the target by $c(\epsilon)=r_\pi(\epsilon)^{-1}$.
The Chinese remainder product $c\,r_\pi$ modulo
$(3)\mathfrak p$ is trivial on all global units, so descends
to the ray class group. If (B4.2) is nontrivial, $c$ has
order three and is nontrivial on the principal units
modulo $(1-\omega)^2$; it cannot factor through the
two-element unit group modulo $1-\omega$. Its conductor
at that prime is exactly the square. At $\mathfrak p$,
$r_\pi$ is nontrivial and has conductor $\mathfrak p$.
This proves the exact finite conductor.

For completeness, the infinity-type alternative is also
well-defined, not just a necessary unit test. Choose
$k$ with $r_\pi(\epsilon)=\epsilon^k$ for every unit.
On principal ideals prime to $\mathfrak p$ define
$r_\pi(\alpha)^{-1}(\alpha/|\alpha|)^k$ for a generator
$\alpha$. Changing the generator to $\epsilon\alpha$
does not change this value. Its finite local component on
residue units is $r_\pi$: multiplying the principal-ideal
value by $r_\pi(\alpha)$ and by the infinity component
$(\alpha/|\alpha|)^{-k}$ gives one. Thus the resulting
ideal character, with this infinity component and local
conductor, is a unitary Hecke character. Its angular
component is nontrivial when $k\equiv2$ or $4\pmod6$,
so it has infinite order by connectedness as above.
\end{proof}

This corrects the finite-order clause in the updated first-pass
note, \S2b item~10. Neutralization at infinity and finite-order
neutralization at the prime over three are different constructions.
For $p=269$, $N(p)=269^2\equiv1\pmod9$, consistent with
the conductor $(269)$ proved in AC-K04.

\subsection{The divisor-pair statement and its positivity limit}

\begin{proposition}[AC-B05: exact norm congruences and Fourier counting]
\label{AC-B05}
Let $R\ge2$ and let $d,e$ be integers prime to $R$. Then
\[
\begin{split}
 N_K(d+e\omega)&=d^2-de+e^2,\\
 N_K(d-e\omega)&=d^2+de+e^2.
\end{split}                                                  \tag{B5.1}
\]
The two respective divisibility tests by $R$ are equivalent
to $d/e$ being a root modulo $R$ of
$\Phi_6(T)=T^2-T+1$ or $\Phi_3(T)=T^2+T+1$.
For either polynomial, roots exist precisely when
\[
 v_3(R)\le1,\qquad
 \ell\equiv1\pmod3\quad\hbox{for every prime }\ell\mid R,
                                                 \ \ell\ne3. \tag{B5.2}
\]
When they exist, their number is
$2^{\#\{\ell\mid R:\ell\ne3\}}$.

For any finite multiset $S$ of integers prime to $R$ and
$\zeta\in(\mathbb Z/R\mathbb Z)^\times$, the ordered-pair
count is exactly
\[
 \#\{(d,e)\in S^2:d\equiv\zeta e\pmod R\}
 =\frac1{\varphi(R)}\sum_{\chi\bmod R}
          \overline{\chi(\zeta)}
          \left|\sum_{d\in S}\chi(d)\right|^2.              \tag{B5.3}
\]
Summing (B5.3) over the roots of the appropriate polynomial
gives the exact norm-divisor-pair count. It does not establish
that a prescribed Erd\H{o}s--Straus divisor set occupies any
of these congruence classes, and it does not give positivity
of the phase-weighted form on all complex test vectors.
\end{proposition}
\begin{proof}
Multiply each element of (B5.1) by its conjugate, using
$\omega+\omega^2=-1$ and $\omega^3=1$. Division by the
unit $e^2$ gives the polynomial criteria. For
$\ell\ne2,3$, the roots of $\Phi_3$ are the two primitive
cube roots in $\mathbb F_\ell^\times$, so exist exactly
when $3\mid\ell-1$. The roots of $\Phi_6$ are their
negatives. In either case the roots are simple and lift
uniquely to every $\ell^a$ by the elementary simple-root
lifting argument. There are no roots at two. At three
there is one root, namely $1$ for $\Phi_3$ and $-1$ for
$\Phi_6$. Substitution $1+3u$ or $-1+3u$ gives a value
of exact 3-adic valuation one, so no root exists modulo
nine. The Chinese remainder theorem proves (B5.2) and
the count. At a composite modulus the local choices can
be made independently; they need not come from one
preselected choice of a global Eisenstein phase.

Character orthogonality on $(\mathbb Z/R\mathbb Z)^\times$
gives
$\mathbf1_{d=\zeta e}=\varphi(R)^{-1}
\sum_\chi\chi(d)\overline{\chi(\zeta e)}$.
Summing over the two copies of $S$ proves (B5.3), including
multiplicities. The distinct roots specify disjoint
congruence conditions for $d/e$, so summing gives the
norm condition exactly.

There is a direct obstruction to promoting this counting
identity to arbitrary-vector positivity. On the cyclic
three-point set let $U$ be cyclic translation. The form
for the union of its two nonidentity phases is represented
by $U+U^*$, whose matrix has zeros on the diagonal and
ones elsewhere. Its eigenvalues are $2,-1,-1$; the vector
$(1,\omega,\omega^2)$ has eigenvalue $-1$. Its values on
nonnegative counting vectors are nonnegative, but it is
not a positive Hermitian form. Finally, if either $d$ or
$e$ is not a unit, the division argument is invalid;
for instance $d=e=R$ satisfies both norm divisibilities
without specifying any unit ratio. These qualifications
exhibit the precise limits of the proposed phase bridge.
\end{proof}

\subsection{The common cyclotomic field and the surviving cubic factors}

\begin{proposition}[AC-B06: the biquadratic base change is exact but not a closure]
\label{AC-B06}
Put $M=\mathbb Q(\zeta_{12})=\mathbb Q(i,\omega)$.
On its Galois group the three nontrivial quadratic
characters satisfy $\chi_{-4}\chi_{-3}\chi_{12}=1$, and
\[
 \zeta_M(s)=\zeta(s)L(s,\chi_{-4})L(s,\chi_{-3})L(s,\chi_{12}).
                                                               \tag{B6.1}
\]
Thus A proves RH for $\zeta_M$. Moreover
\[
 N_{M/\mathbb Q}(r-s\zeta_{12})=r^4-r^2s^2+s^4.                \tag{B6.2}
\]
For a rational prime $p$, let $E'=M(\sqrt[3]p)$ and
let $\eta$ be the quadratic character of $M/K$, namely
the restriction of $\chi_{-4}$, or equivalently its
norm pullback to $K$. Then
\[
\begin{split}
 \operatorname{Gal}(E'/\mathbb Q)&\simeq S_3\times C_2,\\
 L_M(s,\chi_{p,3}|_{G_M})
       &=L_K(s,\chi_{p,3})L_K(s,\chi_{p,3}\eta),\\
 \zeta_{E'}(s)&=\zeta_M(s)
          L_K(s,\chi_{p,3})^2L_K(s,\chi_{p,3}\eta)^2.
\end{split}                                                \tag{B6.3}
\]
All equalities include ramified Euler factors. Consequently
the nonabelian factors have not been eliminated by this
cyclotomic base change or by the product of the three signs.
Under A, RH for $\zeta_{E'}$ is equivalent to RH for both
cubic factors displayed in (B6.3).
\end{proposition}
\begin{proof}
The unit classes $1,5,7,11$ modulo twelve give the table
\[
\begin{array}{c|rrrr}
 &1&5&7&11\\\hline
 \chi_{-4}&1&1&-1&-1\\
 \chi_{-3}&1&-1&1&-1\\
 \chi_{12}&1&-1&-1&1.
\end{array}
\]
The product assertion and (B6.1) follow from this table
and abelian Artin factorization. The product is an
identity of Galois characters, or of characters on
units, not the constant-one identity for their
zero-extended values at all integers. Formula (B6.2)
is $s^4\Phi_{12}(r/s)$, with the case $s=0$ obtained
by the polynomial identity.

The splitting field $E_p$ of $X^3-p$ has group $S_3$
and unique quadratic subfield $K$. Hence
$E_p\cap\mathbb Q(i)=\mathbb Q$, and the compositum
$E'=E_p\mathbb Q(i)$ has group $S_3\times C_2$.
Restriction followed by induction for the index-two
subgroup $G_M\subset G_K$ gives
\[
 \operatorname{Ind}_{G_M}^{G_K}
      (\chi_{p,3}|_{G_M})
              =\chi_{p,3}\oplus\chi_{p,3}\eta.
\]
Artin induction of local inertia invariants proves
the middle identity in (B6.3) at every prime, not only
the unramified ones. The two nontrivial characters
of the cubic extension $E'/M$ are interchanged by
complex conjugation, which preserves ideal norms
and the quadratic character $\eta$. Their global
$L$-functions are therefore equal. Applying the
cyclic Artin factorization over $M$ proves the last
identity. The Hecke functions there are entire,
because both characters are nontrivial. No poles
can cancel their zeros in the open critical strip.
Together with A for $\zeta_M$, this proves the final
equivalence. The field $E'$ is not CM, since it
contains the non-CM-overfield obstruction $E_p$
of AM-H03; adjoining $i$ has not removed that
obstruction.
\end{proof}

\paragraph{Geometric interpretation and an attribution gate.}
Over $\operatorname{Spec}\mathbb Z[1/6]$, the three
quadratic characters are rank-one lisse Artin sheaves
whose tensor product is the unit object; (B6.2) is the
norm of the finite \'etale cyclotomic algebra. Over
$\operatorname{Spec}\mathcal O_M[1/(3p)]$, the cubic
cover is a finite \'etale $\mu_3$-torsor. Its character
summands exist over a coefficient field containing
$\mu_3$, are locally pure of weight zero, and still
have the global functions in (B6.3). These are precise
uses of Galois categories, Kummer torsors and Artin
motives; local weight zero is not a proof of their
global critical-line assertion. Relevant verified
framework locators are recorded in the source appendix.
The updated note attributes $\Delta=\chi_{12}\circ u$
to ``DELTA23 intake L1''. That intake and the map $u$
have not been located among the supplied files in this
run. The elementary character table above is checked;
the specific Grothendieck attribution remains a
source gate and is not asserted as a verified fact.
